% Generated draft: requires mathematical and source review.
\documentclass[11pt]{article}
\usepackage[margin=25mm]{geometry}
\usepackage{fontspec}
\setmainfont{Latin Modern Roman}
\usepackage{amsmath,amssymb,mathtools,unicode-math}
\setmathfont{Latin Modern Math}
\usepackage{enumitem,needspace,xurl,hyperref}
\setlength{\parindent}{0pt}
\setlength{\parskip}{5pt}
\setlength{\emergencystretch}{3em}
\newcommand{\N}{\mathbb{N}}
\newcommand{\Z}{\mathbb{Z}}
\newcommand{\Q}{\mathbb{Q}}
\newcommand{\R}{\mathbb{R}}
\newcommand{\C}{\mathbb{C}}
\newcommand{\F}{\mathbb{F}}
\newcommand{\A}{\mathbb{A}}
\newcommand{\PP}{\mathbb{P}}
\newcommand{\E}{\mathbb{E}}
\newcommand{\eps}{\varepsilon}
\newcommand{\dd}{\,\mathrm{d}}
\begin{document}
\section*{UnsolvedMath ID: 800710}
\textbf{Source problem code:} catalogue entry
\subsection{Definitions and mathematical statement}
Let \(F\) be a field and let \(A\) be a finite-dimensional central simple \(F\)-algebra of degree \(n\), so that \(\dim_F A=n^2\). An involution of the first kind on \(A\) is an \(F\)-linear anti-automorphism \(\sigma\colon A\to A\) satisfying \(\sigma^2=\operatorname{id}_A\). It is orthogonal if, after extension to an algebraic closure of \(F\), it is adjoint to a nondegenerate symmetric nonalternating bilinear form. It is anisotropic if
\[
\forall x\in A,\qquad \sigma(x)x=0\ \Longrightarrow\ x=0.
\]
Let \(\operatorname{SB}(A)\) be the Severi--Brauer variety of \(A\): for a field extension \(E/F\), its \(E\)-points parametrize right ideals of \(A\otimes_F E\) of reduced dimension \(1\), equivalently of \(E\)-dimension \(n\). Write \(F_A=F(\operatorname{SB}(A))\) for its function field, and write \(\sigma_{F_A}=\sigma\otimes\operatorname{id}_{F_A}\).

The supplied Problem 20 asks: for every such \(F\), \(A\), and anisotropic orthogonal involution \(\sigma\), must \(\sigma_{F_A}\) be anisotropic? Explicitly, must one have
\[
\forall y\in A\otimes_F F_A,\qquad
\sigma_{F_A}(y)y=0\ \Longrightarrow\ y=0?
\]
The supplied record states no restriction on \(\operatorname{char}(F)\). The definitions above make sense also in characteristic \(2\), but whether the original source imposes a surrounding characteristic restriction is unresolved from the supplied excerpt.
\subsection{Short English statement}
Does every anisotropic orthogonal involution on a central simple algebra remain anisotropic over the function field of that algebra's Severi--Brauer variety? The supplied excerpt leaves the intended characteristic scope unqualified.
\subsection{Sources}
\begin{itemize}
\item Local UnsolvedMath registry, supplied complete source record: permanent ID \(800710\), exact source code \texttt{catalogue entry}, title \emph{Open Problems on Central Simple Algebras --- Problem 20}.
\item Auel, Brussel, Garibaldi, and Vishne, \emph{Open Problems on Central Simple Algebras} (2010), Problem 20. Supplied source URL: \url{https://arxiv.org/abs/1006.3304}. The registry attributes the record to these authors and gives proposed year 2010.
\item Supplied provenance: accessed 2026-10-16; extraction method \texttt{arxiv\_latex}. This access date is supplied provenance, not an independently verified source visit. The registry labels the problem \texttt{open}; that label is supplied curation, not an independently verified literature assessment.
\end{itemize}
\end{document}
